\section{Question and result}
\label{sec:question}

Resonant-mass gravitational antennas have long been described in terms of elastic eigenmodes, absorption cross sections, effective coupling, and directivity \cite{Hirakawa1976,PaikWagoner1976,Lobo1995}.  The integrated response is also historical: resonant-bar literature treats frequency-integrated absorption cross sections, and adding a resonant transducer redistributes the response among additional normal modes rather than removing its finite material scaling \cite{Aguiar2010}.  Material-response dispersion relations and gravitational-wave sum rules provide another cumulative-response viewpoint \cite{Srivastava2003}.  Complete generator--receiver calculations are historical as well \cite{Rudenko2003}.

The generic mathematical architecture is equally established.  Orthogonal source--receiver wave channels and limits on their connection strengths are standard wave theory \cite{Miller2000}; two-body radiative bounds combine constrained endpoint response with free-space propagation \cite{Molesky2020}; and passive $H_2$ and infinite-dimensional realization machinery are classical systems theory \cite{BarasBrockett1975,Opmeer2013}.  None of these ingredients is claimed as new.

The narrower question addressed here is whether those ingredients imply a simple cumulative ceiling when \emph{both} gravitational endpoints are allowed arbitrary passive linear-harmonic complexity:

\begin{quote}
Can passive resonant engineering increase the frequency-integrated coherent transfer of a separated compact gravitational link without bound, or is there a matter resource that cuts the link at either endpoint?
\end{quote}

Let $\omega_0$ be the absolute gravitational carrier frequency and $\nu$ the complex-envelope detuning, with bandwidth $B/\omega_0\ll1$.  Let $T(\nu)$ map selected local input channels at source $A$ to selected local output channels at receiver $B$, with all unobserved passive ports retained in the physical model.  Define
\begin{equation}
\boxed{
\Gamma_{\rm coh}
\equiv
\frac{1}{2\pi}
\int_{\mathcal B_\nu}\dd\nu\,
\Tr[T^\dagger(\nu)T(\nu)] .
}
\label{eq:gamma}
\end{equation}
For a scalar channel, $\Gamma_{\rm coh}$ is the area under the power-transmissivity spectrum of the complex envelope.  It has units of inverse time and is not an information capacity.

Define the scalar second mass moment about each endpoint center of mass,
\begin{equation}
I_2\equiv\int_V\rho(\bm x)r^2\dd^3x .
\label{eq:I2}
\end{equation}
This is not the moment of inertia about one selected axis; the trace of the conventional inertia tensor is $2I_2$.

The main result is the following leading narrowband wave-zone bound.

\begin{theorem}[Two-ended passive gravitational throughput]
For separated compact nonrelativistic linear-harmonic source and receiver systems in weak mass-quadrupole gravity, with finite or countably infinite bounded-port Markov modal sectors,
\begin{equation}
\boxed{
\Gamma_{\rm coh}
\lesssim
\frac{25G\omega_0^2}{12c^3R^2}
\min(I_{2,A},I_{2,B}) .
}
\label{eq:headline}
\end{equation}
The symbol $\lesssim$ denotes the retained narrowband, leading separated-wave-zone order.  Repeated passive returns between the same two compact endpoints do not change the leading $1/R^2$ coefficient.
\end{theorem}

The result is classical.  Its content is the disappearance, from the leading integrated ceiling, of explicit passive quality factor, finite or countably infinite bounded-port resonance count, passive unitary internal mode mixing, individual modal effective mass, compact quadrupole orientation after optimization, and same-endpoint passive recurrence.  These quantities can strongly reshape the transfer spectrum, but within the stated class they cannot increase its cumulative area beyond the two endpoint resources and free-space TT channel.
